\documentclass[11pt,openany]{article}

\input{algebra-note-preamble}
\input{tcolorbox}
\input{theorem}
\input{algebra-note-commands}

\setstretch{1.25}
\begin{document}
\pagenumbering{arabic}
\begin{center}
	\huge\textbf{Algebra}\\
	\vspace{0.5em}
	\large{Ji, Yong-hyeon}\\
	\vspace{0.5em}
	\normalsize{\today}\\
\end{center}
\begin{note}
	Let \( F \) be a field, and let \( f(x) \in F[x] \) be a polynomial in the polynomial ring over \( F \). For any \( a \in F \), the \textbf{evaluation homomorphism} \[
	\fullfunction{\phi_a}{F[x]}{F}{f(x)}{f(a)}
	\] is defined by \( \varphi_a(f(x)) = f(a) \), where \( f(a) \) denotes the evaluation of \( f(x) \) at \( x = a \).
\end{note}
\thmbox[The Evaluation Homomorphism Theorem]{\begin{theorem}
	 The evaluation map
	\[
	\phi_a : F[x] \to F, \quad f(x) \mapsto f(a)
	\]
	is a \textbf{ring homomorphism}. That is, \( \phi_a \) satisfies the following properties for all \( f(x), g(x) \in F[x] \): \begin{enumerate}[(1)]
		\item \( \phi_a(f(x) + g(x)) = \phi_a(f(x)) + \varphi_a(g(x)) \),
		\item \( \phi_a(f(x) \cdot g(x)) = \phi_a(f(x)) \cdot \varphi_a(g(x)) \).
	\end{enumerate}
\end{theorem}}
\begin{proof}
	\begin{enumerate}[(1)]
		\item Let \( f(x) = \sum_{i=0}^{n} f_i x^i \) and \( g(x) = \sum_{j=0}^{m} g_j x^j \), where \( f_i, g_j \in F \).
		
		Then, \( f(x) + g(x) = \sum_{i=0}^{\max(n,m)} (f_i + g_i) x^i \) (assuming \( f_i = 0 \) or \( g_i = 0 \) for terms where degrees don’t match).
		
		Now, apply the evaluation homomorphism:
		
		\[
		\varphi_a(f(x) + g(x)) = (f(x) + g(x))\big|_{x=a} = \sum_{i=0}^{\max(n,m)} (f_i + g_i) a^i
		\]
		
		On the other hand, applying the evaluation homomorphism separately to \( f(x) \) and \( g(x) \):
		
		\[
		\varphi_a(f(x)) = f(a) = \sum_{i=0}^{n} f_i a^i, \quad \varphi_a(g(x)) = g(a) = \sum_{i=0}^{m} g_j a^j
		\]
		
		Thus:
		
		\[
		\varphi_a(f(x)) + \varphi_a(g(x)) = \sum_{i=0}^{n} f_i a^i + \sum_{j=0}^{m} g_j a^j = \sum_{i=0}^{\max(n,m)} (f_i + g_i) a^i
		\]
		
		Hence:
		
		\[
		\varphi_a(f(x) + g(x)) = \varphi_a(f(x)) + \varphi_a(g(x))
		\]
		
		Thus, the evaluation map preserves **addition**.
	\end{enumerate}
\end{proof}



%#### 2. **Homomorphism for Multiplication:**

Let \( f(x) = \sum_{i=0}^{n} f_i x^i \) and \( g(x) = \sum_{j=0}^{m} g_j x^j \), where \( f_i, g_j \in F \).

The product \( f(x) \cdot g(x) \) is given by:

\[
f(x) \cdot g(x) = \sum_{k=0}^{n+m} \left( \sum_{i+j=k} f_i g_j \right) x^k
\]

Now, apply the evaluation homomorphism:

\[
\varphi_a(f(x) \cdot g(x)) = (f(x) \cdot g(x))\big|_{x=a} = \sum_{k=0}^{n+m} \left( \sum_{i+j=k} f_i g_j \right) a^k
\]

On the other hand, applying the evaluation homomorphism separately to \( f(x) \) and \( g(x) \):

\[
\varphi_a(f(x)) = f(a) = \sum_{i=0}^{n} f_i a^i, \quad \varphi_a(g(x)) = g(a) = \sum_{j=0}^{m} g_j a^j
\]

Thus:

\[
\varphi_a(f(x)) \cdot \varphi_a(g(x)) = \left( \sum_{i=0}^{n} f_i a^i \right) \cdot \left( \sum_{j=0}^{m} g_j a^j \right) = \sum_{k=0}^{n+m} \left( \sum_{i+j=k} f_i g_j \right) a^k
\]

Hence:

\[
\varphi_a(f(x) \cdot g(x)) = \varphi_a(f(x)) \cdot \varphi_a(g(x))
\]

Thus, the evaluation map preserves **multiplication**.

%#### 3. **Homomorphism for Scalar Multiplication:**

For any scalar \( c \in F \), we have \( \varphi_a(c) = c \). This is because \( \varphi_a(c) = c \cdot 1 = c \), as \( c \) is a constant polynomial.

%#### Conclusion:

Since \( \varphi_a \) preserves addition, multiplication, and scalar multiplication, it is a **ring homomorphism**. Therefore, the evaluation map:

\[
\varphi_a : F[x] \to F
\]

defined by \( \varphi_a(f(x)) = f(a) \), is a homomorphism.

%### Final Conclusion (Evaluation Homomorphism Theorem):

\[
\boxed{\varphi_a(f(x)) = f(a) \text{ is a ring homomorphism from } F[x] \text{ to } F.}
\]

\newpage

Let's break down the abstraction from the evaluation map \( f: \mathbb{C} \to \mathbb{C} \) (which takes a point \( a \in \mathbb{C} \) and returns \( f(a) \)) to the more general map \( f: \langle x - a \mid a \in \mathbb{C} \rangle \to \coprod_a \mathbb{C}[x] / \langle x - a \rangle \) (which maps the ideal \( \langle x - a \rangle \) to the corresponding equivalence class \( \overline{f} \in \mathbb{C}[x] / \langle x - a \rangle \)).
%
%### Step-by-Step Explanation:
%
%#### 1. **Evaluation Map: \( f: \mathbb{C} \to \mathbb{C} \)**
This is the familiar **evaluation homomorphism** for a polynomial \( f(x) \in \mathbb{C}[x] \).

- For each \( a \in \mathbb{C} \), the evaluation homomorphism \( \varphi_a \) is defined as:
\[
\varphi_a(f(x)) = f(a)
\]
where \( f(a) \) is simply the result of plugging \( x = a \) into the polynomial \( f(x) \).

In this case, we are thinking of the map as \( f: \mathbb{C} \to \mathbb{C} \), where for each \( a \in \mathbb{C} \), you get a value \( f(a) \in \mathbb{C} \). This is the evaluation map applied to each point \( a \in \mathbb{C} \).

%### 2. **Abstraction to Ideals \( \langle x - a \rangle \)**
We can reinterpret the evaluation of a polynomial \( f(x) \) at \( a \in \mathbb{C} \) in terms of **quotient rings** and **ideals**.

- The **ideal** \( \langle x - a \rangle \) in \( \mathbb{C}[x] \) consists of all polynomials that vanish at \( x = a \), i.e., those divisible by \( x - a \).
- The **quotient ring** \( \mathbb{C}[x] / \langle x - a \rangle \) consists of equivalence classes of polynomials where two polynomials \( f(x) \) and \( g(x) \) are equivalent if they differ by a multiple of \( x - a \). Formally:
\[
f(x) \equiv g(x) \pmod{x - a} \quad \iff \quad f(a) = g(a).
\]

In particular, there is a natural isomorphism between \( \mathbb{C}[x] / \langle x - a \rangle \) and \( \mathbb{C} \), because evaluating a polynomial at \( a \) is equivalent to reducing it modulo \( x - a \):
\[
\mathbb{C}[x] / \langle x - a \rangle \cong \mathbb{C}.
\]
This is because for any \( f(x) \in \mathbb{C}[x] \), the equivalence class \( \overline{f} \in \mathbb{C}[x] / \langle x - a \rangle \) corresponds to the value \( f(a) \in \mathbb{C} \).

%### 3. **The Coproduct**
The coproduct (or direct sum) notation \( \coprod_a \mathbb{C}[x] / \langle x - a \rangle \) refers to the direct sum of the quotient rings \( \mathbb{C}[x] / \langle x - a \rangle \) over all possible values of \( a \in \mathbb{C} \).

- For each \( a \in \mathbb{C} \), we have a quotient ring \( \mathbb{C}[x] / \langle x - a \rangle \), which captures the equivalence classes of polynomials modulo \( x - a \).
- The coproduct \( \coprod_a \mathbb{C}[x] / \langle x - a \rangle \) can be thought of as the "collection" of these quotient rings for all \( a \in \mathbb{C} \).

In this context, the map \( f \) is being extended to take the ideal \( \langle x - a \rangle \) and map it to the corresponding class \( \overline{f} \) in the quotient ring \( \mathbb{C}[x] / \langle x - a \rangle \).

%### 4. **Formal Definition of the New Map**

We now define the map:
\[
f: \langle x - a \mid a \in \mathbb{C} \rangle \to \coprod_a \mathbb{C}[x] / \langle x - a \rangle
\]

- For each \( a \in \mathbb{C} \), we associate the ideal \( \langle x - a \rangle \) with the quotient ring \( \mathbb{C}[x] / \langle x - a \rangle \).
- The map sends the ideal \( \langle x - a \rangle \) to the equivalence class \( \overline{f} \in \mathbb{C}[x] / \langle x - a \rangle \), where \( \overline{f} \) corresponds to the value \( f(a) \).

Formally:
\[
f(\langle x - a \rangle) = \overline{f} \in \mathbb{C}[x] / \langle x - a \rangle.
\]
Since the quotient ring \( \mathbb{C}[x] / \langle x - a \rangle \cong \mathbb{C} \), the image of \( f \) under this map is \( f(a) \in \mathbb{C} \), the value of the polynomial evaluated at \( a \).

%### Summary of Abstraction Process

- The initial map \( f: \mathbb{C} \to \mathbb{C} \) evaluated the polynomial \( f(x) \) at each \( a \in \mathbb{C} \), returning \( f(a) \in \mathbb{C} \).
- This map is abstracted by associating each point \( a \in \mathbb{C} \) with the ideal \( \langle x - a \rangle \) in \( \mathbb{C}[x] \), which corresponds to all polynomials vanishing at \( a \).
- The quotient ring \( \mathbb{C}[x] / \langle x - a \rangle \) captures the equivalence classes of polynomials modulo \( x - a \), and is isomorphic to \( \mathbb{C} \).
- The map is then extended to send \( \langle x - a \rangle \) to the equivalence class \( \overline{f} \in \mathbb{C}[x] / \langle x - a \rangle \), which corresponds to \( f(a) \in \mathbb{C} \).
- The coproduct \( \coprod_a \mathbb{C}[x] / \langle x - a \rangle \) represents the collection of these quotient rings for all \( a \in \mathbb{C} \).

In summary, this process abstracts the simple idea of evaluating a polynomial at each point \( a \) into a more general framework involving ideals, quotient rings, and coproducts, but the core idea remains that the evaluation at each point \( a \) corresponds to a quotient by the ideal \( \langle x - a \rangle \).

\newpage

Yes, this abstraction can be understood and explained more clearly using **category theory**. In category theory, we often use the concepts of **functors**, **universal properties**, and **colimits (coproducts)** to describe structures like the one you're asking about.

Let’s go through how category theory helps formalize and explain this situation.

%### 1. **Objects and Morphisms** in \( \mathbf{Ring} \)
We begin by identifying the relevant objects and morphisms in the category of rings.

- The **object** \( \mathbb{C}[x] \) is the polynomial ring over \( \mathbb{C} \).
- For each \( a \in \mathbb{C} \), we associate an **ideal** \( \langle x - a \rangle \) in \( \mathbb{C}[x] \), which is an ideal in the category of \( \mathbb{C}[x] \)-modules.
- The quotient ring \( \mathbb{C}[x] / \langle x - a \rangle \) is an object in the category of rings, and each quotient ring is isomorphic to \( \mathbb{C} \).

%### 2. **The Evaluation Map as a Functor**

The **evaluation map** at \( a \), denoted \( \varphi_a: \mathbb{C}[x] \to \mathbb{C} \), is a **ring homomorphism**. It takes a polynomial \( f(x) \in \mathbb{C}[x] \) and evaluates it at \( x = a \), giving \( f(a) \in \mathbb{C} \).

In the language of category theory, this can be understood as a **functor** from the category of \( \mathbb{C}[x] \)-modules to the category of \( \mathbb{C} \)-modules:
\[
\varphi_a: \mathbf{Mod}_{\mathbb{C}[x]} \to \mathbf{Mod}_{\mathbb{C}}.
\]
This functor is **covariant** and sends each polynomial \( f(x) \) to its evaluation \( f(a) \) in \( \mathbb{C} \).

%### 3. **Universal Properties of Quotients**

The quotient ring \( \mathbb{C}[x] / \langle x - a \rangle \) has a **universal property** in category theory. It is the largest quotient ring of \( \mathbb{C}[x] \) that "kills" the ideal \( \langle x - a \rangle \), meaning it maps \( x \) to \( a \) via the natural homomorphism.

For each \( a \in \mathbb{C} \), the quotient map:
\[
q_a: \mathbb{C}[x] \to \mathbb{C}[x] / \langle x - a \rangle \cong \mathbb{C}
\]
is a morphism in the category \( \mathbf{Ring} \), and it satisfies the universal property that every ring homomorphism \( \mathbb{C}[x] \to R \) that sends \( x \) to \( a \) factors uniquely through the quotient ring \( \mathbb{C}[x] / \langle x - a \rangle \).

Thus, this quotient captures the behavior of the evaluation at \( a \).

%### 4. **Coproduct in Category Theory**

The coproduct \( \coprod_a \mathbb{C}[x] / \langle x - a \rangle \) is the **categorical coproduct** (also known as the **direct sum**) of all the quotient rings \( \mathbb{C}[x] / \langle x - a \rangle \) over all \( a \in \mathbb{C} \). In the category \( \mathbf{Ring} \), the coproduct of rings is essentially a **disjoint union** of the quotient rings.

In this case, the coproduct takes the form:
\[
\coprod_{a \in \mathbb{C}} \mathbb{C}[x] / \langle x - a \rangle,
\]
where each component \( \mathbb{C}[x] / \langle x - a \rangle \) is isomorphic to \( \mathbb{C} \). The coproduct provides a way to "combine" all the evaluations of \( f(x) \) at different points \( a \in \mathbb{C} \).

In category theory, the coproduct is the **universal object** that maps into each of the components (the quotient rings in this case) via natural inclusion maps. In other words, for each evaluation map \( \varphi_a \), there is a corresponding projection map from the coproduct to the component \( \mathbb{C}[x] / \langle x - a \rangle \).

%### 5. **The Functorial Interpretation of the Process**

Let’s think of the entire process in terms of **functors**.

- The family of ideals \( \langle x - a \rangle \), indexed by \( a \in \mathbb{C} \), forms a **diagram** in the category \( \mathbf{Ring} \).
- The evaluation map \( f: \langle x - a \rangle \to \mathbb{C}[x] / \langle x - a \rangle \) can be seen as a **natural transformation** between the functor that takes each ideal \( \langle x - a \rangle \) to the corresponding quotient ring \( \mathbb{C}[x] / \langle x - a \rangle \).
- The coproduct \( \coprod_{a \in \mathbb{C}} \mathbb{C}[x] / \langle x - a \rangle \) is the **colimit** (coproduct) of this functor, collecting all possible evaluations into one structure.

%### 6. **Abstracting from Pointwise Evaluation to Ideals**

In category theory, the transition from the pointwise evaluation map \( f: \mathbb{C} \to \mathbb{C} \), where \( f(a) = f(x) \big|_{x=a} \), to the map \( f: \langle x - a \rangle \to \mathbb{C}[x] / \langle x - a \rangle \) is an abstraction. Instead of working directly with points in \( \mathbb{C} \), we now work with **ideals** in the polynomial ring \( \mathbb{C}[x] \).

The ideals \( \langle x - a \rangle \) represent the vanishing conditions at the points \( a \in \mathbb{C} \), and the quotient rings \( \mathbb{C}[x] / \langle x - a \rangle \) encapsulate the idea of evaluating a polynomial at those points. The coproduct then captures the global behavior of the polynomial at all points simultaneously.

%### Conclusion

Using **category theory**, we can interpret the abstraction from pointwise evaluation to quotient rings as a process of defining a functor from the category of ideals to the category of quotient rings, with the coproduct representing the "sum" of all evaluations. The universal properties of the quotient rings and the coproduct ensure that this abstraction is both well-defined and captures the essence of polynomial evaluation in a more general, categorical framework.

In this sense, category theory provides a powerful language for understanding and abstracting the process of evaluating polynomials at points into a more general and structured setting involving ideals, quotient rings, and coproducts.

\newpage

Certainly! Before reinterpreting the process of evaluation homomorphisms, ideals, and quotient rings into category theory, we need to review several fundamental **definitions** and **theorems** from ring theory and category theory. These will form the foundation for understanding the arguments in categorical terms.

%### **Necessary Definitions and Theorems**:

%#### **1. Ring Homomorphisms**
A **ring homomorphism** between two rings \( R \) and \( S \) is a function \( \psi: R \to S \) that preserves the ring operations:
- **Addition**: \( \psi(r_1 + r_2) = \psi(r_1) + \psi(r_2) \).
- **Multiplication**: \( \psi(r_1 \cdot r_2) = \psi(r_1) \cdot \psi(r_2) \).
- **Identity**: \( \psi(1_R) = 1_S \).

%#### **2. Polynomial Ring**
Let \( R \) be a commutative ring. The **polynomial ring** \( R[x] \) consists of all polynomials in the indeterminate \( x \) with coefficients in \( R \). For example:
\[
f(x) = \sum_{i=0}^{n} r_i x^i \quad \text{where} \quad r_i \in R.
\]
- **Addition and multiplication** of polynomials are defined in the usual way, inheriting the ring structure from \( R \).

%#### **3. Evaluation Homomorphism**
The **evaluation homomorphism** \( \varphi_a: R[x] \to R \), for a fixed \( a \in R \), is the map that evaluates the polynomial \( f(x) \in R[x] \) at \( x = a \):
\[
\varphi_a(f(x)) = f(a).
\]
- **Kernel**: The kernel of \( \varphi_a \) is the ideal \( \langle x - a \rangle \), the set of all polynomials that vanish at \( x = a \).

%#### **4. Ideal in a Ring**
An **ideal** \( I \) in a ring \( R \) is a subset \( I \subseteq R \) such that:
1. \( I \) is an additive subgroup of \( R \).
2. For any \( r \in R \) and \( i \in I \), we have \( r \cdot i \in I \) (closed under multiplication with elements from the ring).

In this context, the ideal \( \langle x - a \rangle \) is the principal ideal generated by \( x - a \), which consists of all polynomials in \( R[x] \) that are divisible by \( x - a \).

%#### **5. Quotient Ring**
Given a ring \( R \) and an ideal \( I \subseteq R \), the **quotient ring** \( R/I \) is the set of cosets of \( I \) in \( R \):
\[
R/I = \{ r + I \mid r \in R \}.
\]
The operations of addition and multiplication are defined naturally on cosets. The quotient ring \( R/I \) is a ring, and the projection map \( \pi: R \to R/I \) is a ring homomorphism with kernel \( I \).

%#### **6. First Isomorphism Theorem (Rings)**
Let \( \psi: R \to S \) be a ring homomorphism. The **First Isomorphism Theorem** states that:
\[
R / \ker(\psi) \cong \operatorname{Im}(\psi),
\]
where \( \ker(\psi) \) is the kernel of \( \psi \) and \( \operatorname{Im}(\psi) \) is the image of \( \psi \). This theorem allows us to identify the quotient ring \( R / \ker(\psi) \) with the image of \( \psi \).

%#### **7. Categorical Concepts:**

- **Categories**: A **category** consists of objects and morphisms between them, where morphisms represent structure-preserving maps. For example, the category \( \mathbf{Ring} \) consists of rings as objects and ring homomorphisms as morphisms.

- **Functors**: A **functor** \( F: \mathcal{C} \to \mathcal{D} \) between two categories \( \mathcal{C} \) and \( \mathcal{D} \) is a map that assigns to each object in \( \mathcal{C} \) an object in \( \mathcal{D} \), and to each morphism in \( \mathcal{C} \) a corresponding morphism in \( \mathcal{D} \), while preserving composition and identities.

- **Natural Transformations**: A **natural transformation** is a way of transforming one functor into another, in a manner that is "natural" with respect to the objects and morphisms in the category.

- **Limits and Colimits**: These are universal constructions in category theory. The **coproduct** (or **direct sum**) is a specific type of colimit, representing the "sum" of a collection of objects in a category.

%#### **8. Coproduct (Direct Sum) in Ring Theory**
The **coproduct** in the category of rings, denoted \( \coprod \), is a type of colimit that generalizes the notion of the **direct sum**. In the context of rings, the coproduct of a family of rings \( \{R_i\} \) is the "freest" ring that embeds each \( R_i \). For quotient rings \( R_i = \mathbb{C}[x] / \langle x - a_i \rangle \), we have:
\[
\coprod_{a_i} \mathbb{C}[x] / \langle x - a_i \rangle.
\]

This coproduct corresponds to a disjoint union of the quotient rings.

---

%### **Reinterpreting in Category Theory:**

Now that we've established the basic definitions and theorems, let’s explain the key ideas that arise from them in category theory.

1. **Evaluation Homomorphism as a Functor:**
- The evaluation homomorphism \( \varphi_a: \mathbb{C}[x] \to \mathbb{C} \) can be interpreted as a functor from the category of \( \mathbb{C}[x] \)-modules to \( \mathbb{C} \)-modules. This functor evaluates polynomials at specific points \( a \in \mathbb{C} \).

2. **Quotient Ring and Isomorphism:**
- By the **First Isomorphism Theorem**, we know that \( \mathbb{C}[x] / \langle x - a \rangle \cong \mathbb{C} \), so the evaluation map \( \varphi_a \) factors through the quotient ring:
\[
\mathbb{C}[x] \to \mathbb{C}[x] / \langle x - a \rangle \cong \mathbb{C}.
\]

3. **Coproduct as a Colimit:**
- The coproduct \( \coprod_a \mathbb{C}[x] / \langle x - a \rangle \) is a colimit in the category of rings. It represents the "direct sum" of all possible evaluations of polynomials at different points \( a \in \mathbb{C} \). Each component of the coproduct corresponds to a different quotient ring, and the coproduct encapsulates the global structure of evaluating a polynomial at every point in \( \mathbb{C} \).

4. **Abstraction to Functors:**
- The map \( f: \mathbb{C} \to \mathbb{C} \), where \( f(a) = f(x) \big|_{x=a} \), is now abstracted as a functor from the category of ideals \( \langle x - a \rangle \) to the category of quotient rings \( \mathbb{C}[x] / \langle x - a \rangle \).
- The evaluation map \( f: \langle x - a \rangle \to \mathbb{C}[x] / \langle x - a \rangle \) can be interpreted as a natural transformation between functors, associating each ideal \( \langle x - a \rangle \) with its corresponding quotient ring.

---

%### Summary:

By combining these definitions and theorems with category theory, we reinterpret the **evaluation map** as a **functor** and the **quotient rings** as objects in a coproduct (a type of colimit). This categorical perspective provides a deeper understanding of how polynomial evaluation, quotient rings, and coproducts are connected in the category of rings.

In short:
- **Evaluation homomorphism** becomes a **functor** from ideals to quotient rings.
- **Quotient rings** \( \mathbb{C}[x] / \langle x - a \rangle \) are the objects associated with each point \( a \).
- The **coproduct** (direct sum) of these quotient rings captures the global structure of evaluations at all points.

\newpage

Yes, it is possible to state a formal proposition summarizing the abstraction you are asking about. Here’s how we can state it:

---

**Proposition:**

Let \( \mathbb{C}[x] \) be the ring of polynomials over \( \mathbb{C} \), and for each \( a \in \mathbb{C} \), let \( \langle x - a \rangle \) be the principal ideal of \( \mathbb{C}[x] \) generated by \( x - a \). There exists a natural homomorphism:

\[
f: \{\langle x - a \rangle \mid a \in \mathbb{C} \} \to \coprod_{a \in \mathbb{C}} \mathbb{C}[x] / \langle x - a \rangle
\]
such that for each \( a \in \mathbb{C} \), the ideal \( \langle x - a \rangle \) is mapped to the equivalence class \( \overline{f} \in \mathbb{C}[x] / \langle x - a \rangle \), where \( \overline{f} \) is the image of \( f(x) \) in the quotient ring \( \mathbb{C}[x] / \langle x - a \rangle \).

Furthermore, for each \( a \in \mathbb{C} \), the quotient ring \( \mathbb{C}[x] / \langle x - a \rangle \) is isomorphic to \( \mathbb{C} \), and the map \( f \) can be interpreted as an evaluation map that assigns to each ideal \( \langle x - a \rangle \) the evaluation of the polynomial \( f(x) \) at \( a \), i.e.,

\[
f(\langle x - a \rangle) = f(a) \in \mathbb{C}.
\]

By the **First Isomorphism Theorem**, we have:

\[
\mathbb{C}[x] / \langle x - a \rangle \cong \mathbb{C}.
\]

Thus, the coproduct:

\[
\coprod_{a \in \mathbb{C}} \mathbb{C}[x] / \langle x - a \rangle
\]

collects the evaluations of the polynomial \( f(x) \) at all points \( a \in \mathbb{C} \).

---

**Explanation of the Proposition**:

- **Objects**: The ideal \( \langle x - a \rangle \) represents polynomials vanishing at \( x = a \). The quotient ring \( \mathbb{C}[x] / \langle x - a \rangle \) represents the evaluation of polynomials at \( a \), and this quotient is isomorphic to \( \mathbb{C} \).

- **Map \( f \)**: The map \( f \) takes each ideal \( \langle x - a \rangle \) and sends it to the corresponding evaluation of the polynomial \( f(x) \) at \( x = a \), represented as the equivalence class \( \overline{f} \in \mathbb{C}[x] / \langle x - a \rangle \).

- **Coproduct**: The coproduct \( \coprod_{a \in \mathbb{C}} \mathbb{C}[x] / \langle x - a \rangle \) collects the evaluations of \( f(x) \) at every point in \( \mathbb{C} \), forming the direct sum of the quotient rings.

---

Why this Proposition is Valid:

- **First Isomorphism Theorem**: The key part of this proposition is the First Isomorphism Theorem, which tells us that the quotient \( \mathbb{C}[x] / \langle x - a \rangle \) is isomorphic to \( \mathbb{C} \), making the evaluation at \( a \) correspond to the quotient by \( \langle x - a \rangle \).

- **Coproduct**: The coproduct collects all these quotient evaluations for each \( a \in \mathbb{C} \), providing a categorical way to manage evaluations at all points in \( \mathbb{C} \).

---

This proposition provides a precise and formal statement about the structure you're asking about, combining ideas from polynomial rings, quotient rings, evaluation maps, and coproducts.

\newpage

You are right to point out the distinction between interpreting a polynomial \( f(x) \in \mathbb{C}[x] \) as a function from \( \mathbb{C} \to \mathbb{C} \) by **evaluation** (pointwise evaluation at a complex number), and the more abstract interpretation where we use **ideals**, **quotient rings**, and possibly a **coproduct** structure. Let's go through this step by step to clarify the distinction.

1. **Interpretation of \( f \in \mathbb{C}[x] \) as a Function \( f: \mathbb{C} \to \mathbb{C} \)**

In the classical sense, a polynomial \( f(x) \in \mathbb{C}[x] \) can be viewed as a function:
\[
f: \mathbb{C} \to \mathbb{C}, \quad f(a) = f(x)\big|_{x=a}.
\]

This is the **evaluation map** at \( x = a \). In this case:

- **Domain**: \( \mathbb{C} \), where each element \( a \in \mathbb{C} \) is a point at which we evaluate the polynomial.
- **Codomain**: \( \mathbb{C} \), which is the set of complex numbers, representing the result of the evaluation \( f(a) \in \mathbb{C} \).

Here, both the domain and codomain are simply \( \mathbb{C} \), and **evaluation** means that for each \( a \in \mathbb{C} \), we compute \( f(a) \) using the polynomial.

2. **Using the First Isomorphism Theorem**

Now, when we involve **ideals** and **quotient rings**, we're looking at a different perspective.

- The **evaluation map** \( \varphi_a: \mathbb{C}[x] \to \mathbb{C} \) takes the polynomial \( f(x) \in \mathbb{C}[x] \) and evaluates it at \( x = a \).

- The **kernel** of this evaluation map \( \varphi_a \) is the ideal \( \langle x - a \rangle \), which consists of all polynomials in \( \mathbb{C}[x] \) that vanish at \( x = a \).

The **First Isomorphism Theorem** tells us that the quotient ring \( \mathbb{C}[x] / \langle x - a \rangle \) is isomorphic to \( \mathbb{C} \). That is:
\[
\mathbb{C}[x] / \langle x - a \rangle \cong \mathbb{C}.
\]

This is important because it identifies the **structure** of the quotient ring, but it doesn’t necessarily describe the basic **pointwise evaluation** of polynomials in \( \mathbb{C}[x] \). The First Isomorphism Theorem is used to relate the kernel of the evaluation map (the ideal \( \langle x - a \rangle \)) to the quotient structure. This is why the **domain** of the map, which is based on these ideals, uses the First Isomorphism Theorem.

Why the Domain Uses the First Isomorphism Theorem:
- The First Isomorphism Theorem helps identify that for each \( a \in \mathbb{C} \), the quotient \( \mathbb{C}[x] / \langle x - a \rangle \) is isomorphic to \( \mathbb{C} \), making the quotient ring the "structure" corresponding to evaluation at \( x = a \).
- So when we treat \( f(x) \) as a map from \( \mathbb{C} \to \mathbb{C} \), the **First Isomorphism Theorem** is relevant because it tells us that the quotient ring structure captures the evaluation process at each point \( a \).

3. **Why the Codomain Doesn’t Use a Coproduct in Pointwise Evaluation**

When we are simply thinking of \( f: \mathbb{C} \to \mathbb{C} \) by pointwise evaluation, there's no need to introduce a coproduct because:

- **Pointwise evaluation** is a simple function: for each \( a \in \mathbb{C} \), the map gives \( f(a) \in \mathbb{C} \). 
- In this context, there's no need to consider combining evaluations at multiple points at the same time (which is what the coproduct would do). Instead, we're just looking at evaluation at a **single** point \( a \), so the codomain is simply \( \mathbb{C} \).

Why the Codomain Doesn’t Use a Coproduct for Pointwise Evaluation:
- In the pointwise evaluation setting, we’re simply mapping each point \( a \in \mathbb{C} \) to \( f(a) \in \mathbb{C} \). Since there’s no need to consider the evaluations at multiple points at once, the codomain is just \( \mathbb{C} \), not a coproduct.

4. **When the Codomain Uses a Coproduct**

When we move to the more abstract setting involving **ideals** and **quotient rings**, and we consider **all** evaluations of \( f(x) \) at every point \( a \in \mathbb{C} \), then the coproduct structure becomes relevant.

- The **coproduct** \( \coprod_{a \in \mathbb{C}} \mathbb{C}[x] / \langle x - a \rangle \) collects the evaluations of \( f(x) \) at all points simultaneously.
- Each quotient ring \( \mathbb{C}[x] / \langle x - a \rangle \cong \mathbb{C} \) represents the evaluation at a single point \( a \), and the coproduct organizes all these evaluations together.

In the coproduct setting, we are no longer looking at **individual evaluations** at one point but instead at **evaluations at all points** \( a \in \mathbb{C} \) simultaneously, and that’s why the codomain in this case is a coproduct.

Summary:

- When \( f(x) \in \mathbb{C}[x] \) is thought of as a function \( f: \mathbb{C} \to \mathbb{C} \), we are doing **pointwise evaluation** at a single point. In this case:
- The **domain** is \( \mathbb{C} \), which consists of the points \( a \) where we evaluate.
- The **codomain** is \( \mathbb{C} \), where the result of the evaluation lives.
- There's no need for a coproduct in the codomain because we are evaluating at **one point** at a time.

- When we abstract the process to consider **ideals** and **quotient rings**:
- The **First Isomorphism Theorem** helps us describe the structure of the quotient ring \( \mathbb{C}[x] / \langle x - a \rangle \cong \mathbb{C} \), which is why the **domain** involves the First Isomorphism Theorem.
- The **coproduct** in the codomain collects the evaluations of \( f(x) \) at all points \( a \in \mathbb{C} \), combining the results of evaluation across all points simultaneously.

\end{document}
